\section{DIS is diverse}
\label{sec:diverse}
% \begin{itemize}
%     \item $\gamma/Z$
%     \item leptons
%     \item hadrons: p+d+A; nucl. PDF
%     \item experimental overview
% \end{itemize}

In this section we give an overview many different DIS datasets commonly used in PDF fits.
By doing so, we quickly realize that the picture in \cref{chap:intro} is far too simplistic.

\subsection{HERA}
We start with the Hadron-Elektron-Ringanlage (HERA),  which was in operation between 1992 and 2007.
It was hosted by DESY in Hamburg, Germany.
\fherror{add HERA image}
While the promoted physics discovery, leptoquarks, eventually did not happen, HERA has become one of the most influential DIS machines so far.
In fact, it is until this day the only lepton-proton collider in the world.
It recorded a huge amount of data, which is still analyzed nowadays, 20 years after the shutdown.
HERA had four experiments (H1, ZEUS, HERA-B and HERMES), but we focus here on the combined analysis of H1 and ZEUS~\cite{H1:2015ubc}.
In \cref{tab:HERA} we given an overview of the most important fully-inclusive DIS measurements from HERA.

\begin{table}
    \caption{HERA measurement~\cite{H1:2015ubc} used in NNPDF4.0~\cite{NNPDF:2021njg}}
    \label{tab:HERA}
    \begin{tabular}{ll ll ll}
        projectile & $\sqrt s$ [GeV] & x & $Q^2$ & $N_\text{dat}^\text{fit}/N_\text{dat}$ & obs.\\
        \hline
        $e^-$ & 318 & ? & ? & ?/201-43 & $\sigma_r^\text{NC}$ \\
        $e^+$ & 318 & ? & ? & ? & $\sigma_r^\text{NC}$ \\
        $e^+$ & 300 & ? & ? & ? & $\sigma_r^\text{NC}$ \\
        $e^+$ & 251 & ? & ? & ? & $\sigma_r^\text{NC}$ \\
        $e^+$ & 225 & ? & ? & ? & $\sigma_r^\text{NC}$ \\
        $e^-$ & 318 & ? & ? & ? & $\sigma_r^\text{CC}$ \\
        $e^+$ & 318 & ? & ? & ? & $\sigma_r^\text{CC}$ 
    \end{tabular}
\end{table}
\fherror{complete HERA table}

We start by looking at the projectile column and immediately we realize that HERA was most of the time not colliding electrons, but positrons.
This is to remind us that DIS is lepton-hadron scattering and, indeed, we have to consider all leptons.
We expand our selection further below in \cref{sec:neutrino}, but remain for the moment with the electron and the positron, which are exact copies of each other, but with opposite charge.

Next, we turn to the energy column $\sqrt s$, which spreads across four different energies.
In practice the lepton beam was operated with a fixed beam energy ($E_e = \SI{27.5}{\GeV}$) and varying energies for the protons ($E_p=\SIlist{920;820;575;460}{\GeV}$).
If we then look to the $x$ and $Q^2$ columns an recall \cref{eq:defshad}, we see that the table only indicates the maximal available range of those variables.
Specifically, HERA was able to measure activity in the detector for $0.005 < y < 0.95$.

However, not only the accessible kinematical phase space limits our kinematic coverage, but also our requirement to study DIS, \cref{box:DISreq}.
In practice, e.g.\ NNPDF4.0~\cite{NNPDF:2021njg} applies a cut of $Q^2 \geq \SI{8}{\GeV^2}$ and $W^2 \geq \SI{12.5}{\GeV^2}$.
Mathematically, the $Q^2$ cut ensures we are in a perturbative regime, \cref{eq:CpQCD}, while the $W^2$ cut has multiple motivations.
First, it ensures we are in an inelastic regime, i.e.\ we are not finding resonances or excitations of the proton.
Second and related to the former, we know it suppresses the correction term $\order{\frac{\LQCD^2}{Q^2}}$ in \cref{eq:fact1}.
Studying and trying to estimate the correction term is a topic in current research~\cite{Ball:2025xtj,Harland-Lang:2025wvm,Cerutti:2025yji}.
Looking to the column with the number of data points, we note that even after the kinematic cuts we are left with a total of \fherror{$N_\text{data}^\text{HERA}$} measured cross sections.

Finally, we turn to the observables column for which we see two different entries exist:
the reduced neutral current (NC) cross section $\sigma_r^\text{NC}$ and the reduced charge current (CC) cross section $\sigma_r^\text{CC}$.
The \enquote{current} here refers to the current between the leptonic part and the hadronic part, in particular the exchanged boson.
The \enquote{reduced} cross section refers to a slight rescaling of the double differential cross section, \cref{eq:dsigmadxdy},
\begin{equation}
    \sigma_r = \frac{xy Q^2}{4\pi \aem^2}\frac{\dd[2] \sigma}{\dd x \dd y}
\end{equation}
such that it refers to a specific PDF combination at LO accuracy.

Next, we examine the CC case by recalling two important facts from \cref{box:particles}:
first, charged leptons, and thus also electrons and positrons, have weak charges and, second, $W$ bosons have an electric charge.
Thus, by CC we consider the lepton interacts with the hadron via the exchange of a $W$ boson, which behaves quite different then a photon.
First and foremost, due to the electric charge conservation the $W$ boson can only interact with specific leptons and quarks.

\subsection{Neutrino DIS}
\label{sec:neutrino}
